\setcounter{Exercise}{0}

\chapter*{Feuille d'exercices : Les variables}

\section*{Exemple :}

Soit le code suivant :
\begin{lstlisting}[style=pythonstyle, language=Python]
PI = 3.1415
rayon = float(input("Saisissez le rayon :"))
perimetre = 2 * PI * rayon
print("Le perimetre du cercle est", perimetre, "m")  

\end{lstlisting}

\subsection*{Question 1}
Dans ce programme, combien y a-t-il de variables ? Quels sont leur identifiant ?

\subsection*{Question 2}
Donner ligne par ligne dans un tableau la valeur des variables.
On supposera que l'utilisateur a saisi un rayon de 2 m.\\

\subsection*{Question 3}
Expliquez-en un paragraphe ce que fait le programme.

\vfill

\section*{Exercice formulaire :}

Soit le code suivant :
\begin{lstlisting}[style=pythonstyle, language=Python]
nom = input("Saisissez votre nom :")
prenom = input("Saisissez votre prenom :")
age = int(input("Saisissez votre age :"))

print("Vous vous appelez", prenom, nom, "et vous avez", age, "ans.")  
\end{lstlisting}

\subsection*{Question 1}
Dans ce programme, combien y a-t-il de variables ? Quels sont leur identifiant ?

\subsection*{Question 2}
Donner ligne par ligne dans un tableau la valeur des variables.
On supposera que l'utilisateur s'appelle John Doe et a 35 ans.\\

\subsection*{Question 3}
Expliquez-en un paragraphe ce que fait le programme.

\vfill

\section*{Exercice multiplication :}
\begin{lstlisting}[style=pythonstyle, language=Python]
a = int(input("Saisissez un nombre :"))
b = int(input("Saisissez un nombre :"))
res = a * b 
print("Le produit de", a, "et", b, "vaut", res)  
\end{lstlisting}

\subsection*{Question 1}
Dans ce programme, combien y a-t-il de variables ? Quels sont leur identifiant ?

\subsection*{Question 2}
Donner ligne par ligne dans un tableau la valeur des variables.
On supposera que l'utilisateur a saisi les nombres 2 et 24.\\

\subsection*{Question 3}
Expliquez-en un paragraphe ce que fait le programme.

\vfill

\section*{Exercice chichoumeille : Niveau 1}
\begin{lstlisting}[style=pythonstyle, language=Python]
a = 3
b = a - 1
c = b * a
a = 12
a = 2 * a - (c * b)
\end{lstlisting}
\subsection*{Question 1}
Dans ce programme, combien y a-t-il de variables ? Quels sont leur identifiant ?

\subsection*{Question 2}
Donner ligne par ligne dans un tableau la valeur des variables.


\vfill

\section*{Exercice chichoumeille : Niveau 2}
\begin{lstlisting}[style=pythonstyle, language=Python]
a = 3
b = 6
c = b // a
c = 2**c
b = c % a
\end{lstlisting}
\subsection*{Question 1}
Dans ce programme, combien y a-t-il de variables ? Quels sont leur identifiant ?

\subsection*{Question 2}
Donner ligne par ligne dans un tableau la valeur des variables.

\vfill

\section*{Exercice chichoumeille : Niveau 3}
\begin{lstlisting}[style=pythonstyle, language=Python]
a = 1
b = a > 0
a = (a+1)**2
b = a <= 4
b = (a % 2) == 0
\end{lstlisting}
\subsection*{Question 1}
Dans ce programme, combien y a-t-il de variables ? Quels sont leur identifiant et leur type ?

\subsection*{Question 2}
Donner ligne par ligne dans un tableau la valeur des variables.


\section*{Exercice chichoumeille : Hardcore}
\begin{lstlisting}[style=pythonstyle, language=Python]
a = 15
b = (a % 3 + 1)**a
c = (b * 5 // 3) >= 2
a = a + a / 2
b = b % (b+1) - a
\end{lstlisting}
\subsection*{Question 1}
Dans ce programme, combien y a-t-il de variables ? Quels sont leur identifiant ?

\subsection*{Question 2}
Donner ligne par ligne dans un tableau la valeur des variables.